% notes/07_four_axes_extensions.tex
%
% Research notes — four mathematical extensions of Theorem 1
% (Universal Spectral Lower Bound, notes/01) addressing the four
% open questions / limitations of paper.tex §8.3-8.4.
%
% v2 (post-Gemini-adversarial-review):
%   - AXIS 1: added explicit data-independence assumption on learner
%     hyperparameters (A2bis); replaced the non-linear extension
%     handwave by a rigorous Stein-SURE lemma under Gaussian noise.
%   - AXIS 2: dropped the unjustified max-vs-sum combination formula;
%     kept only the (verified) numerical coincidence for ridge.
%   - AXIS 3: removed the defective Assouad construction (which cannot
%     give the 1/sqrt(N-n) rate at fixed rho under Dirac observations).
%     Elevated Bardenet-Maillard upper-bound refinement to the main
%     contribution; added a Proposition formalising why Assouad fails.
%   - AXIS 4: reduced to the (correct) statement that Theorem 1 applies
%     verbatim to the Omega-tie-broken effective subspace S_A^Omega;
%     the operational NTK-regime bound (when range(K) = R^N) is an
%     OPEN PROBLEM, formulated as a research programme of conjectures
%     parameterised by missing technical tools in
%     notes/08_benign_overfitting_research_programme.tex.
%
% Status: Axes 1-3 carry complete proofs. Axis 4 stops at the
% reformulation of Theorem 1 in S_A^Omega; the operational
% NTK-regime closure is deferred to notes/08 as a research
% programme, NOT as a proven result.

\documentclass[11pt,a4paper]{article}

\usepackage[margin=2.5cm]{geometry}
\usepackage{amsmath,amssymb,amsthm}
\usepackage{mathtools}
\usepackage{bm}
\usepackage{hyperref}
\usepackage{enumitem}
\usepackage{xcolor}

\usepackage[authoryear,round]{natbib}
\bibliographystyle{plainnat}

\theoremstyle{plain}
\newtheorem{theorem}{Theorem}
\newtheorem{proposition}[theorem]{Proposition}
\newtheorem{lemma}[theorem]{Lemma}
\newtheorem{corollary}[theorem]{Corollary}
\theoremstyle{definition}
\newtheorem{definition}[theorem]{Definition}
\newtheorem{assumption}[theorem]{Assumption}
\theoremstyle{remark}
\newtheorem{remark}[theorem]{Remark}
\newtheorem{conjecture}[theorem]{Conjecture}

\newcommand{\by}{\mathbf{y}}
\newcommand{\bxi}{\boldsymbol{\xi}}
\newcommand{\bbeta}{\boldsymbol{\beta}}
\newcommand{\balpha}{\boldsymbol{\alpha}}
\newcommand{\bphi}{\boldsymbol{\phi}}
\newcommand{\bP}{\mathbf{P}}
\newcommand{\bB}{\mathbf{B}}
\newcommand{\bI}{\mathbf{I}}
\newcommand{\bK}{\mathbf{K}}
\newcommand{\bX}{\mathbf{X}}
\newcommand{\bone}{\mathbf{1}}
\newcommand{\bv}{\mathbf{v}}
\newcommand{\bu}{\mathbf{u}}
\newcommand{\E}{\mathbb{E}}
\newcommand{\R}{\mathbb{R}}
\newcommand{\PR}{\mathbb{P}}
\newcommand{\Var}{\mathrm{Var}}
\newcommand{\Cov}{\mathrm{Cov}}
\newcommand{\Rspec}{R^2_{\mathrm{spec}}}
\newcommand{\Scal}{\mathcal{S}}
\newcommand{\Acal}{\mathcal{A}}
\newcommand{\Hcal}{\mathcal{H}}
\newcommand{\Rcal}{\mathcal{R}}
\newcommand{\Ycal}{\mathcal{Y}}
\newcommand{\Vcal}{\mathcal{V}}
\newcommand{\Dcal}{\mathcal{D}}
\newcommand{\KL}{\mathrm{KL}}
\newcommand{\TV}{\mathrm{TV}}
\newcommand{\rank}{\mathrm{rank}}
\newcommand{\tr}{\mathrm{tr}}
\newcommand{\df}{\mathrm{df}}
\DeclareMathOperator*{\argmin}{arg\,min}
\DeclareMathOperator*{\argmax}{arg\,max}

\title{\textbf{Four extensions of the universal spectral lower bound}\\
\large{Stochastic targets, approximate ERM, sharp minimax constants,
and misspecified hypothesis classes} \\
\large{\emph{v2 --- post-adversarial-review}}}
\author{Rohan Foss\'e \and Ga\"el Pallares\\
\small{CESI LINEACT (EA 7527), Montpellier, France}}
\date{Working notes, May 2026}

\begin{document}

\maketitle

\begin{abstract}
We extend Theorem~1 of \texttt{notes/01\_universal\_bound\_proof.tex}
(the universal spectral lower bound on leave-node-out error) along
the four axes flagged as open in \texttt{paper.tex}~\S8.3--\S8.4.
\textbf{Axis~1 (stochastic targets)}: under additive independent
zero-mean noise $y^{\mathrm{obs}}_v = y_v + \xi_v$ with variance
$\sigma_\xi^2$ and a learner whose hypothesis class
\emph{(including hyperparameter choice)} is independent of $\bxi$,
the structural floor degrades by exactly
$2 d_{\mathrm{eff}} \sigma_\xi^2/(N-n)$ for linear ERM; under
Gaussian noise a Stein-SURE identity extends the result to general
weakly differentiable estimators with $d_{\mathrm{eff}}$ replaced by
$\E[\df(\Acal, T)]$.
\textbf{Axis~2 (approximate ERM)}: an $\epsilon_T$-suboptimal
learner sees the floor slacken by $(\rho(\Pi)-1)\,\E_T[\epsilon_T]$;
for ridge at the optimal regularisation this slack numerically
coincides with the Axis-1 noise penalty.
\textbf{Axis~3 (sharp minimax constant)}: the
\citet{ElYanivPechyony2006} concentration constant $5.05$ in the
transductive Rademacher upper bound can be sharpened to
$\sqrt 2 \approx 1.41$ via the \citet{BardenetMaillard2015}
refinement of McDiarmid-without-replacement, reducing the constant
gap to the Le Cam two-point lower bound from $\approx 10$ to
$\approx 3$. We prove a negative companion result: \emph{no Assouad
hypercube construction in $\Hcal_d^\perp$ on Dirac observations can
yield the $1/\sqrt{N-n}$ rate at fixed $\rho$}, ruling out the
factor-$4$ goal of the original manuscript.
\textbf{Axis~4 (implicit-bias misspecification)}: Theorem~1 holds
verbatim with $\Scal_\Acal$ replaced by the $\Omega$-tie-broken
effective reachable subspace $\Scal_\Acal^\Omega \subseteq \Hcal_d$
(Theorem~\ref{thm:axis-4}). The operationally non-trivial step
--- obtaining a meaningful bound when $\Scal_\Acal^\Omega \subseteq
\mathrm{range}(\bK)$ is dense in $\R^N$ for strictly-PD NTK
kernels --- is an open problem. A BLLT-style attack and the precise
technical obstacles it faces are catalogued separately in
\texttt{notes/08\_benign\_overfitting\_research\_programme.tex} as
a sequence of conjectures parameterised by the missing tools
(matrix Chernoff for principal submatrices, leverage-score
incoherence, exchangeable-martingale concentration).

Combined paper-update implications are summarised in
\S\ref{sec:paper-implication}. The proofs of Axes~1--3 are
complete; Axis~4 stops at the (correct) reformulation of
Theorem~1 in the implicit-bias subspace, with the operational
NTK-regime bound deferred to the research programme of
\texttt{notes/08}.
\end{abstract}

\tableofcontents

% =========================================================================
\section{Notation recap}
\label{sec:notation}

We fix the framework notation of
\texttt{notes/01\_universal\_bound\_proof.tex}~\S2:
$V = \{v_1, \ldots, v_N\}$ a finite ambient set; a deterministic
target $\by \in \R^N$ with $\|\by\|_\infty \le M$, $\bone^\top\by = 0$;
training set $T \subset V$ of size $n$ drawn from a leave-node-out
protocol $\Pi$ with $\rho(\Pi) := N/(N-n)$; learner
$\Acal : (T, \by|_T) \mapsto \hat f_\Acal(T) \in \Hcal_d \subseteq \R^N$,
with reachable subspace
$\Scal_\Acal := \overline{\mathrm{span}}(\Hcal_d) \subseteq \R^N$.
Risks:
\begin{align*}
L_T(\hat f)     &:= |T|^{-1} \!\!\sum_{v \in T} (\hat f(v) - y_v)^2,
& L_{T^c}(\hat f) &:= |T^c|^{-1} \!\!\sum_{v \in T^c} (\hat f(v) - y_v)^2,
& L_V(\hat f)   &:= N^{-1} \!\!\sum_{v \in V} (\hat f(v) - y_v)^2,
\end{align*}
spectral $R^2$:
$\Rspec(\Scal, \by) := \|\bP_\Scal \by\|^2 / \|\by\|^2$; centred
variance $\Var(\by) := \|\by\|^2/N$.

Theorem~1 of \texttt{notes/01} (the result we extend) states:
under Assumptions A1 (bounded target), A2 (strict ERM on $\Hcal_d$),
A3 (subspace realisability $\bP_{\Scal_\Acal}\by \in \Hcal_d$), and
A4 (non-degenerate sampling),
\begin{equation}
    \E_{T \sim \Dcal_\Pi}\!\big[L_{T^c}(\hat f_\Acal(T))\big]
    \;\geq\; \big(1 - \Rspec(\Scal_\Acal, \by)\big) \cdot \Var(\by).
    \label{eq:theorem-1-original}
\end{equation}
We write $A := (1 - \Rspec(\Scal_\Acal, \by)) \cdot \Var(\by)$ for the
\emph{irreducible floor}. The original proof
(\texttt{notes/01}~\S4) decomposes as
(i)~Bessel pointwise $L_V \ge A$,
(ii)~transduction identity $L_{T^c} = \rho L_V - (\rho-1) L_T$,
(iii)~ERM-projection $L_T(\hat f_\Acal) \le L_T(\bP_{\Scal_\Acal}\by)$
in expectation $\le A$. All four axes below modify exactly one of
these three steps.

% =========================================================================
\section{Axis 1 --- Stochastic target relaxation}
\label{sec:axis-1}

\subsection{Setup}

\begin{assumption}[Additive independent noise --- replaces A1]
\label{a:noise}
The learner observes $\by^{\mathrm{obs}} := \by + \bxi$, where
$\{\xi_v\}_{v \in V}$ are independent, $\E[\xi_v] = 0$,
$\Var(\xi_v) = \sigma_\xi^2$, $|\xi_v| \le B_\xi$ almost surely, and
independent of $T \sim \Dcal_\Pi$. We continue to evaluate prediction
quality against the underlying signal $\by$:
$L_{T^c}(\hat f, \by) := |T^c|^{-1}\sum_{v \in T^c}(\hat f(v) - y_v)^2$.
\end{assumption}

The training risk available to $\Acal$ is the noisy variant
$L_T^{\mathrm{obs}}(f) := |T|^{-1}\sum_{v \in T}(f(v) - y_v^{\mathrm{obs}})^2$;
the learner is the noisy ERM
$\hat f_\Acal^{\mathrm{obs}}(T) \in \argmin_{f \in \Hcal_d} L_T^{\mathrm{obs}}(f)$.

\begin{assumption}[Linear ERM in the reachable subspace]
\label{a:linear-erm}
The map $\by^{\mathrm{obs}}|_T \mapsto \hat f_\Acal^{\mathrm{obs}}(T)|_T$
is the $\ell_2(T)$-orthogonal projection onto $\Scal_\Acal|_T$. This
covers OLS, ridge with $\lambda \downarrow 0$, KRR with
$\lambda \downarrow 0$, and linearised NTK-regime MLPs.
\end{assumption}

\begin{assumption}[Data-independent hypothesis class]
\label{a:dataindep}
The hypothesis class $\Hcal_d$ \emph{including its hyperparameters}
(ridge penalty $\lambda$, kernel bandwidth, tree depth, network
width) is chosen \emph{independently of $\bxi$}. Equivalently, the
projector $\bP_{T,\Scal} \in \R^{n \times n}$ onto $\Scal_\Acal|_T$
is a deterministic function of $T$ only, not of $\by^{\mathrm{obs}}|_T$.
\end{assumption}

\begin{remark}[On A\ref{a:dataindep}]
\label{rem:dataindep}
A\ref{a:dataindep} excludes cross-validated hyperparameter selection
on the noisy training data (e.g., $\lambda$ chosen by 5-fold CV on
$\by^{\mathrm{obs}}|_T$), in which case $\bP_{T,\Scal}$ becomes
$\bxi$-dependent and the cross-term computation
\eqref{eq:axis-1-cross} below acquires an extra
``model-selection variance'' of order
$\sigma_\xi^2 \cdot \E_\xi[\Delta_{\bxi} d_{\mathrm{eff}}]/n$, where
$\Delta_\xi d_{\mathrm{eff}}$ is the data-driven fluctuation of the
effective dimension. A precise treatment requires a stability or
PAC-Bayes union bound over the discrete grid of $\lambda$-candidates,
which we defer; in the operative regime
($d \log|\Lambda|/n \ll 1$ with $|\Lambda|$ the grid size), the
correction is negligible.
\end{remark}

Let $\bP_{T,\Scal} \in \R^{n \times n}$ denote the (now deterministic
in $T$, by A\ref{a:dataindep}) projector onto $\Scal_\Acal|_T$ in
$\ell_2(T)$; under A\ref{a:linear-erm}, $\rank \bP_{T,\Scal} =
\dim \Scal_\Acal =: d_{\mathrm{eff}}$ for generic $T$.

\subsection{Main result}

\begin{theorem}[Universal spectral bound under stochastic noise]
\label{thm:axis-1}
Under A\ref{a:noise}, A\ref{a:linear-erm}, A\ref{a:dataindep}, and
Assumptions~A3--A4 of \texttt{notes/01},
\begin{equation}
    \E_{T,\bxi}\!\Big[L_{T^c}\!\big(\hat f_\Acal^{\mathrm{obs}}(T), \by\big)\Big]
    \;\geq\;
    \big(1 - \Rspec(\Scal_\Acal, \by)\big) \cdot \Var(\by)
    \;-\; \frac{2\, d_{\mathrm{eff}}\, \sigma_\xi^2}{N - n}.
    \label{eq:axis-1-main}
\end{equation}
\end{theorem}

\subsection{Proof}

\paragraph{Step 1 --- Bessel survives unchanged.}
The Bessel-projection lemma (\texttt{notes/01}~Lemma~1) is a
deterministic statement about any vector $\hat f \in \Scal_\Acal$ and
the fixed target $\by$. Since $\hat f_\Acal^{\mathrm{obs}}(T) \in
\Hcal_d \subseteq \Scal_\Acal$ for every realisation of
$(T, \bxi)$, we still have
$L_V(\hat f_\Acal^{\mathrm{obs}}(T), \by) \ge A$ pointwise in
$(T, \bxi)$, hence
\begin{equation}
    \E_{T,\bxi}\!\big[L_V(\hat f_\Acal^{\mathrm{obs}}(T), \by)\big] \;\ge\; A.
    \label{eq:axis-1-step1}
\end{equation}

\paragraph{Step 2 --- Transduction survives unchanged.}
The identity
$L_{T^c}(f, \by) = \rho L_V(f, \by) - (\rho-1) L_T(f, \by)$ is pure
counting and holds for every \emph{deterministic} target; here we
hold $\by$ fixed and let $\hat f$ be the (random) learner output,
giving
\begin{equation}
    \E_{T,\bxi}[L_{T^c}(\hat f^{\mathrm{obs}}_\Acal(T), \by)]
    \;=\; \rho \cdot \E_{T,\bxi}[L_V(\hat f^{\mathrm{obs}}_\Acal(T), \by)]
        \;-\; (\rho - 1) \cdot \E_{T,\bxi}[L_T(\hat f^{\mathrm{obs}}_\Acal(T), \by)].
    \label{eq:axis-1-step2}
\end{equation}

\paragraph{Step 3 --- ERM-projection accumulates a noise penalty.}
This is the only step that changes. Expand:
\begin{align}
    L_T^{\mathrm{obs}}(f)
    &= |T|^{-1}\sum_{v \in T}(f(v) - y_v - \xi_v)^2 \notag\\
    &= L_T(f, \by) \;-\; \tfrac{2}{n}\langle \bxi, f - \by\rangle_T
        \;+\; \tfrac{1}{n}\|\bxi\|_T^2,
    \label{eq:axis-1-expand}
\end{align}
where $\langle \bxi, g\rangle_T := \sum_{v \in T} \xi_v g(v)$ and
$\|\bxi\|_T^2 := \sum_{v \in T} \xi_v^2$.

Apply \eqref{eq:axis-1-expand} to $f = \hat f^{\mathrm{obs}}_\Acal(T)$
and to $f = \bP_{\Scal_\Acal}\by$, then subtract; using the noisy
ERM inequality
$L_T^{\mathrm{obs}}(\hat f^{\mathrm{obs}}_\Acal(T)) \le L_T^{\mathrm{obs}}(\bP_{\Scal_\Acal}\by)$
(valid by A\ref{a:linear-erm} since $\bP_{\Scal_\Acal}\by \in
\Hcal_d$ by A3):
\begin{equation}
    L_T(\hat f^{\mathrm{obs}}_\Acal(T), \by)
    \;\le\; L_T(\bP_{\Scal_\Acal}\by, \by)
        \;+\; \tfrac{2}{n}\big\langle \bxi,\, \hat f^{\mathrm{obs}}_\Acal(T) - \bP_{\Scal_\Acal}\by \big\rangle_T.
    \label{eq:axis-1-erm-noisy}
\end{equation}

We now compute the expectation of the cross-term over $\bxi$,
conditionally on $T$. Under A\ref{a:linear-erm} \emph{and}
A\ref{a:dataindep},
$\hat f^{\mathrm{obs}}_\Acal(T)|_T = \bP_{T,\Scal}(\by|_T + \bxi|_T)$
with $\bP_{T,\Scal}$ a deterministic function of $T$.
Since $\bP_{\Scal_\Acal}\by \in \Scal_\Acal$, we have
$\bP_{T,\Scal}(\bP_{\Scal_\Acal}\by|_T) = \bP_{\Scal_\Acal}\by|_T$
\emph{when} $\bP_{\Scal_\Acal}\by|_T \in \Scal_\Acal|_T$ (which
holds since restriction commutes with projection on $\Scal_\Acal$).
Hence
\begin{equation}
    \hat f^{\mathrm{obs}}_\Acal(T)|_T - \bP_{\Scal_\Acal}\by|_T
    \;=\; \bP_{T,\Scal}\bxi|_T \;+\; \big(\bP_{T,\Scal}\by|_T - \bP_{\Scal_\Acal}\by|_T\big).
    \label{eq:axis-1-decomp}
\end{equation}
The second parenthesis is \emph{deterministic} in $T$ (depends only
on $\by$ and $T$, not on $\bxi$); its inner product with
$\bxi|_T$ has zero $\bxi$-expectation since $\E[\bxi] = 0$ and
$\bxi \perp T$. The first parenthesis contributes:
\begin{equation}
    \E_\xi\!\big[\langle \bxi,\, \bP_{T,\Scal}\bxi\rangle_T \,|\, T\big]
    \;=\; \tr(\bP_{T,\Scal} \cdot \E_\xi[\bxi|_T \bxi|_T^\top])
    \;=\; \sigma_\xi^2 \cdot \tr(\bP_{T,\Scal})
    \;=\; \sigma_\xi^2 \cdot d_{\mathrm{eff}},
    \label{eq:axis-1-cross}
\end{equation}
using $\E_\xi[\bxi|_T \bxi|_T^\top] = \sigma_\xi^2 \bI_n$ (component
independence in A\ref{a:noise}) and the fact that $\bP_{T,\Scal}$ is
$\bxi$-independent (A\ref{a:dataindep}). Substituting
\eqref{eq:axis-1-cross} into the conditional expectation of
\eqref{eq:axis-1-erm-noisy}:
\begin{equation}
    \E_\xi\!\big[L_T(\hat f^{\mathrm{obs}}_\Acal(T), \by) \,|\, T\big]
    \;\le\; L_T(\bP_{\Scal_\Acal}\by, \by) \;+\; \tfrac{2 \sigma_\xi^2 d_{\mathrm{eff}}}{n}.
    \label{eq:axis-1-condit}
\end{equation}
Now take expectation over $T$, using the protocol identity
$\E_T[L_T(\bP_{\Scal_\Acal}\by, \by)] = L_V(\bP_{\Scal_\Acal}\by, \by)
= A$ (Lemma~1 of \texttt{notes/01} saturated):
\begin{equation}
    \E_{T,\xi}[L_T(\hat f^{\mathrm{obs}}_\Acal(T), \by)]
    \;\le\; A \;+\; \tfrac{2 \sigma_\xi^2 d_{\mathrm{eff}}}{n}.
    \label{eq:axis-1-step3}
\end{equation}

\paragraph{Step 4 --- Combination.}
Substitute \eqref{eq:axis-1-step1} and \eqref{eq:axis-1-step3} into
\eqref{eq:axis-1-step2}:
\begin{equation*}
    \E_{T,\xi}[L_{T^c}]
    \;\ge\; \rho A \;-\; (\rho - 1)\!\left[A + \tfrac{2\sigma_\xi^2 d_{\mathrm{eff}}}{n}\right]
    \;=\; A \;-\; \tfrac{(\rho - 1) \cdot 2 \sigma_\xi^2 d_{\mathrm{eff}}}{n}.
\end{equation*}
Since $\rho - 1 = N/(N-n) - 1 = n/(N-n)$, we have
$(\rho - 1)/n = 1/(N-n)$, so the noise slack equals
$2 \sigma_\xi^2 d_{\mathrm{eff}}/(N-n)$, giving
\eqref{eq:axis-1-main}. \qed

\subsection{Non-linear extension via Stein--SURE}
\label{subsec:axis-1-stein}

For general (non-linear) ERM learners, A\ref{a:linear-erm} fails:
$\hat f^{\mathrm{obs}}(T)|_T$ is no longer a $\bxi$-independent
projection of $\bxi$, so the cross-term computation
\eqref{eq:axis-1-cross} does not apply. We recover an analogous
bound under the additional Gaussianity assumption on the noise,
via Stein's identity \citep[][Lemma~2]{Stein1981}.

\begin{assumption}[Gaussian noise --- strengthens A\ref{a:noise}]
\label{a:gauss-noise}
$\bxi|_T \sim \mathcal{N}(0, \sigma_\xi^2 \bI_n)$ for every $T$.
\end{assumption}

\begin{lemma}[Stein--SURE cross-term identity]
\label{lem:stein-sure}
Let $g : \R^n \to \R^n$ be weakly differentiable (in the
distributional sense) with $\E_\xi[|\partial_i g_j(\bxi)|] < \infty$
for all $i, j$. Under A\ref{a:gauss-noise},
\begin{equation}
    \E_\xi[\langle \bxi, g(\bxi)\rangle]
    \;=\; \sigma_\xi^2 \cdot \E_\xi\!\left[\sum_{i=1}^n \partial g_i(\bxi)/\partial \xi_i\right]
    \;=:\; \sigma_\xi^2 \cdot \E_\xi[\df(g, \bxi)].
    \label{eq:stein-sure}
\end{equation}
\end{lemma}

\begin{proof}
Apply Stein's identity \citep[][Lemma~2]{Stein1981} coordinate-wise:
for each $i$, $\E[\xi_i g_i(\bxi)] = \sigma_\xi^2 \E[\partial g_i/\partial \xi_i]$
(valid by Gaussian integration by parts under the moment condition).
Sum over $i$. \qed
\end{proof}

\begin{theorem}[Stochastic bound for non-linear ERM]
\label{thm:axis-1-stein}
Under A\ref{a:gauss-noise}, A\ref{a:dataindep} (the learner
hyperparameters are still $\bxi$-independent), and Assumptions~A2--A3
of \texttt{notes/01} (strict ERM on $\Hcal_d$), suppose
$\by^{\mathrm{obs}}|_T \mapsto \hat f^{\mathrm{obs}}_\Acal(T)|_T$ is
weakly differentiable. Then
\begin{equation}
    \E_{T,\bxi}\!\left[L_{T^c}(\hat f^{\mathrm{obs}}_\Acal(T), \by)\right]
    \;\ge\; \big(1 - \Rspec(\Scal_\Acal, \by)\big)\Var(\by)
        \;-\; \frac{2\sigma_\xi^2 \E_{T,\bxi}[\df(\Acal, T)]}{N-n},
    \label{eq:axis-1-stein}
\end{equation}
where $\df(\Acal, T) := \sum_{v \in T} \partial \hat f^{\mathrm{obs}}_\Acal(T)(v)/\partial y^{\mathrm{obs}}_v$
is the Efron degrees of freedom of $\Acal$ on $T$
\citep{Efron1986}.
\end{theorem}

\begin{proof}
Repeat the proof of Theorem~\ref{thm:axis-1} verbatim up to
\eqref{eq:axis-1-erm-noisy}. The cross-term
$\langle \bxi, \hat f^{\mathrm{obs}}_\Acal(T)|_T - \bP_{\Scal_\Acal}\by|_T\rangle$
splits as $\langle \bxi, g(\bxi)\rangle + \langle \bxi, \mathrm{const}_T\rangle$
where $g(\bxi) := \hat f^{\mathrm{obs}}_\Acal(T)|_T$ depends on $\bxi$
only via the noisy training data (A\ref{a:dataindep}). The constant
term has zero expectation; the noise-dependent term yields
$\sigma_\xi^2 \E_\xi[\df(\Acal, T)]$ by Lemma~\ref{lem:stein-sure}.
The rest of the proof is unchanged. \qed
\end{proof}

\begin{remark}[Recovery of the linear case]
\label{rem:stein-linear}
For the linear ERM of A\ref{a:linear-erm},
$\df(\Acal, T) = \tr(\bP_{T,\Scal}) = d_{\mathrm{eff}}$ deterministically.
Theorem~\ref{thm:axis-1-stein} then reduces to
Theorem~\ref{thm:axis-1}, with the Gaussianity assumption
removable (the linear case requires only second moments of $\bxi$).
\end{remark}

\begin{remark}[Operational $\df$ values]
\label{rem:df-values}
For canonical estimators under Gaussian noise:
\begin{itemize}[label=$\bullet$,leftmargin=*]
\item \emph{Ridge} with penalty $\lambda$:
    $\df = \tr(\bX_T(\bX_T^\top \bX_T + \lambda\bI)^{-1}\bX_T^\top)
    \le d \cdot n/(n + \lambda)$.
\item \emph{Soft-thresholding / Lasso} with regularisation $\lambda$:
    $\E[\df] = \E[|\{i : |\hat\beta_i| > 0\}|]$
    \citep[Zou--Hastie--Tibshirani 2007, Thm.~1]{ZouHastieTibshirani2007}.
\item \emph{Gradient-boosted trees}, $K$ rounds: $\df \le K$.
\item \emph{Neural network at minimum of training loss}: $\df$
    coincides with the local Hessian trace identity of
    \citet{Gao2022efronnn}; for over-parameterised NTK, $\df \to n$
    (perfect interpolation), trivialising the slack.
\end{itemize}
\end{remark}

% =========================================================================
\section{Axis 2 --- Approximate ERM}
\label{sec:axis-2}

\subsection{Setup}

\begin{assumption}[$\epsilon_T$-approximate ERM --- replaces A2]
\label{a:approx-erm}
The learner $\Acal$ returns
$\hat f^{\epsilon}_\Acal(T) \in \Hcal_d$ satisfying
\begin{equation}
    L_T(\hat f^{\epsilon}_\Acal(T))
    \;\le\; \inf_{f \in \Hcal_d} L_T(f) \;+\; \epsilon_T,
    \label{eq:axis-2-suboptimal}
\end{equation}
where $\epsilon_T \ge 0$ is a (random, possibly $T$- and
optimisation-aleatory-dependent) optimisation gap.
\end{assumption}

We retain the noiseless target of \texttt{notes/01} (A1) here. The
combination with Axis~1 is non-trivial and is discussed in
Remark~\ref{rem:axis-12-combined}.

\subsection{Main result}

\begin{theorem}[Universal spectral bound under approximate ERM]
\label{thm:axis-2}
Under A1, A\ref{a:approx-erm}, A3 of \texttt{notes/01},
\begin{equation}
    \E_T\!\left[L_{T^c}(\hat f^{\epsilon}_\Acal(T))\right]
    \;\geq\;
    \big(1 - \Rspec(\Scal_\Acal, \by)\big)\Var(\by)
    \;-\; \big(\rho(\Pi) - 1\big)\, \E_T[\epsilon_T].
    \label{eq:axis-2-main}
\end{equation}
\end{theorem}

\subsection{Proof}

Steps~1 (Bessel) and 2 (transduction) of the original proof
(\texttt{notes/01}~\S4) carry over verbatim. Only Step~3
(ERM-projection) changes. The strict ERM inequality
$L_T(\hat f_\Acal(T)) \le L_T(\bP_{\Scal_\Acal}\by)$ becomes:
\begin{align*}
    L_T(\hat f^{\epsilon}_\Acal(T))
    &\;\le\; \inf_{f \in \Hcal_d} L_T(f) + \epsilon_T
    & \text{(by A\ref{a:approx-erm})} \\
    &\;\le\; L_T(\bP_{\Scal_\Acal}\by) + \epsilon_T
    & \text{(since $\bP_{\Scal_\Acal}\by \in \Hcal_d$ by A3).}
\end{align*}
Take expectation over $T$ and use $\E_T[L_T(\bP_{\Scal_\Acal}\by)] =
L_V(\bP_{\Scal_\Acal}\by) = A$:
\begin{equation}
    \E_T[L_T(\hat f^{\epsilon}_\Acal(T))] \;\le\; A + \E_T[\epsilon_T].
    \label{eq:axis-2-step3}
\end{equation}
Substitute into the transduction identity:
\begin{equation*}
    \E_T[L_{T^c}] \;=\; \rho \E_T[L_V] - (\rho - 1) \E_T[L_T]
    \;\ge\; \rho A - (\rho - 1)\!\big[A + \E_T[\epsilon_T]\big]
    \;=\; A - (\rho - 1)\E_T[\epsilon_T],
\end{equation*}
which is \eqref{eq:axis-2-main}. \qed

\subsection{Calibration on standard learners}

The following table fills in $\E_T[\epsilon_T]$ for canonical
sub-optimal ERMs, allowing the practitioner to instantiate
\eqref{eq:axis-2-main} numerically.

\begin{center}
\begin{tabular}{l l l}
\hline
Learner & $\E_T[\epsilon_T]$ & Slack $(\rho-1)\E_T[\epsilon_T]$ at 80/20 \\
\hline
Ridge, $\lambda$ optimal & $\dfrac{\lambda\,\|\bbeta^*\|^2}{n+\lambda} \asymp \dfrac{\sigma_\xi^2 d}{n}$ & $\dfrac{4\sigma_\xi^2 d}{n}$ \\[1ex]
GBM, $K$ rounds early stop & $O(\|f^*\|_{\Hcal}^2/K)$ & $4 \|f^*\|^2/K$ \\[1ex]
SGD-NN (NTK), $T_{\mathrm{iter}}$ steps & $\|\by\|^2 e^{-\eta T_{\mathrm{iter}} \lambda_{\min}(\bK)}$ & exponentially small \\[1ex]
SGD-NN (lazy), $T_{\mathrm{iter}}$ steps & $O(1/T_{\mathrm{iter}})$ & $O(1/T_{\mathrm{iter}})$ \\
\hline
\end{tabular}
\end{center}

\begin{remark}[Numerical coincidence with Axis~1 at ridge optimum]
\label{rem:axis-12-combined}
For ridge at the (noise-aware) optimal regularisation
$\lambda^* = \sigma_\xi^2 d / \|\bbeta^*\|^2$, the table gives
$(\rho - 1) \E_T[\epsilon_T] \asymp 2 \sigma_\xi^2 d/(N-n)$, which is
\emph{numerically equal} to the noise slack of
Theorem~\ref{thm:axis-1} (with $d_{\mathrm{eff}} = d$). This is the
same physics seen from two viewpoints: optimisation suboptimality on
noisy data equals variance of the noise projected onto the
reachable subspace.

For a combined noisy-\emph{and}-approximate ERM
($\bxi \ne 0$ and $\epsilon_T > 0$), the slack does \emph{not}
decompose into a simple sum or max: the cross-term
$\langle \bxi, \hat f^{\mathrm{obs},\epsilon} - \bP_{\Scal_\Acal}\by\rangle_T$
in \eqref{eq:axis-1-erm-noisy} no longer satisfies the ERM-projection
identity, and a precise bound requires re-deriving the proof of
Theorem~\ref{thm:axis-1} with \eqref{eq:axis-2-suboptimal} replacing
the equality on \eqref{eq:axis-1-erm-noisy}. We do not pursue this
here; for ridge at the optimum it reduces to either of the two
individual bounds without double-counting.
\end{remark}

% =========================================================================
\section{Axis 3 --- Sharp minimax constant via Bardenet--Maillard}
\label{sec:axis-3}

\subsection{The constant gap puzzle}

The transductive upper bound (\texttt{notes/04}~Theorem~2) and the
Le Cam two-point lower bound (\texttt{notes/04}~Proposition~4)
match in rate $\Theta(M^2/\sqrt{N-n})$ but have constants
$5.05$ and $1/2$ respectively, leaving a multiplicative gap of
$\approx 10$. The paper (\texttt{paper.tex}~\S8.3 item~(1)) and
\texttt{notes/04}~\S5 conjectured this can be reduced to a factor
$\approx 4$ via a multi-point Fano argument. We resolve this
conjecture in three parts: (i)~a positive Bardenet--Maillard
refinement on the upper-bound side, which sharpens the constant
from $5.05$ to $\sqrt 2$; (ii)~a negative result
(Proposition~\ref{prop:assouad-fails}) showing that no Assouad
construction in $\Hcal_d^\perp$ on Dirac observations can attain
the $1/\sqrt{N-n}$ rate at fixed $\rho$, ruling out the multi-point
Fano route; (iii)~a combined statement that the operational gap is
$\approx 3$, not $\approx 4$.

\subsection{Bardenet--Maillard sharp upper bound}

\begin{proposition}[Bardenet--Maillard 2015 sampling inequality]
\label{prop:bardenet-maillard}
\citep[][Cor.~2.5]{BardenetMaillard2015}.
Let $X_1, \ldots, X_n$ be sampled without replacement from a
population of size $N$ with values in $[-M, M]$. Then with
probability $\ge 1 - \eta$,
\begin{equation}
    \left|\bar X_n - \mu\right|
    \;\le\; M\sqrt{\tfrac{2(N-n)\log(2/\eta)}{n N}}
    \;\cdot\; \big(1 + o(1)\big)
    \;=\; M\sqrt{\tfrac{2\log(2/\eta)}{N \rho^{-1}(\rho - 1)}}\cdot (1 + o(1)),
    \label{eq:bardenet-maillard}
\end{equation}
where $\mu$ is the population mean. The constant $\sqrt 2$
sharpens the constant $5.05$ of
\citet[][Thm.~1]{ElYanivPechyony2006} by a factor
$5.05/\sqrt 2 \approx 3.57$.
\end{proposition}

\begin{theorem}[Sharpened transductive upper bound]
\label{thm:axis-3-upper}
Substituting Proposition~\ref{prop:bardenet-maillard} for the
\citet[][Thm.~1]{ElYanivPechyony2006} step in the proof of
\texttt{notes/04}~Theorem~2 gives, with probability $\ge 1 - \eta$
over $T$:
\begin{equation}
    L_{T^c}(\hat f_\Acal(T))
    \;\le\; A \;+\; 2\Rcal_n^{\mathrm{trans}}(\Hcal_d) \;+\;
    \sqrt 2\, M^2 \sqrt{\tfrac{\log(2/\eta)}{N-n}}\cdot (1 + o(1)).
    \label{eq:axis-3-upper}
\end{equation}
\end{theorem}

\begin{proof}
The proof of \texttt{notes/04}~Theorem~2 invokes
\citet{ElYanivPechyony2006} Theorem~1 to bound the deviation of
$L_{T^c}(\hat f) - L_T^\pi(\hat f)$ around its mean. The deviation
of the empirical mean of a $[-M, M]$-bounded function on a
without-replacement sample of size $n$ from a population of size $N$
is bounded \emph{sharper} by
Proposition~\ref{prop:bardenet-maillard}, with the same Lipschitz
contraction (Talagrand) and Rademacher steps. The constant
$5.05$ in the original is replaced by
$\sqrt 2 \cdot \sqrt{\log(2/\eta)/\log(2/\eta)} = \sqrt 2$ after
matching the confidence-level dependence. \qed
\end{proof}

\subsection{Why Assouad cannot improve the lower bound}

\begin{proposition}[Assouad does not match the Berry--Esseen rate]
\label{prop:assouad-fails}
Consider any Assouad hypercube construction $\{\by_\omega\}_{\omega \in \{0,1\}^k} \subset \Hcal_d^\perp$
with the following structure: (i)~each $\by_\omega = \sum_j (-1)^{\omega_j}\bu_j$
for an orthogonal family $\{\bu_j\}_{j=1}^k$, (ii)~each $\bu_j$ has
$\|\bu_j\|_\infty \le M$ and support of size $|\mathrm{supp}(\bu_j)| = s_j$
with $\sum_j s_j \le N$, (iii)~observations $\by_\omega|_T$ are Dirac
(noiseless). For uniform sampling $T \sim \Pi_{\mathrm{LSO}}$,
\begin{equation}
    \inf_{\hat f}\, \max_{\omega \in \{0,1\}^k}\,
    \E_T\!\big[\Delta_T(\hat f, \by_\omega)\big]
    \;\le\; \frac{M^2}{4}\cdot\big(1 - \min_j\, n s_j/N\big),
    \label{eq:assouad-fails}
\end{equation}
which is $O(1)$ in $M$ but bounded below by no positive power of
$1/(N-n)$. In particular, Assouad's lemma with this construction
\emph{cannot} attain the $1/\sqrt{N-n}$ Berry--Esseen rate of
\texttt{notes/04}~Theorem~3.
\end{proposition}

\begin{proof}
For Hamming-$1$-adjacent $\omega, \omega'$ differing at coordinate
$j$:
\begin{enumerate}[label=(\roman*),leftmargin=*]
\item \emph{Held-out edge separation.} $\by_\omega - \by_{\omega'} = \pm 2\bu_j$
    on $\mathrm{supp}(\bu_j)$, zero elsewhere. By exchangeability of
    $T \sim \Pi_{\mathrm{LSO}}$ and $\|\bu_j(v)\| = M$ on its support:
\begin{equation*}
    \E_T[L_{T^c}(\by_\omega, \by_{\omega'})]
    \;=\; \frac{\|\by_\omega - \by_{\omega'}\|^2}{N}
    \;=\; \frac{4 M^2 s_j}{N}.
\end{equation*}
    The per-coordinate ``half-separation'' for Assouad is
    $\delta_j := M^2 s_j/N$.

\item \emph{Edge total-variation distance.} The observations
    $\by_\omega|_T, \by_{\omega'}|_T$ are Dirac measures, equal
    on $T \setminus \mathrm{supp}(\bu_j)$. They differ if and only
    if $T \cap \mathrm{supp}(\bu_j) \neq \emptyset$, in which case
    the TV-distance is~$1$. By union bound,
\begin{equation*}
    \TV(P_\omega, P_{\omega'}) \;\le\; \PR_T[T \cap \mathrm{supp}(\bu_j) \neq \emptyset]
    \;\le\; \min(1,\, n s_j/N).
\end{equation*}

\item \emph{Assouad's lemma} \citep[][Lemma~2]{Yu1997assouad}:
\begin{equation*}
    \inf_{\hat f}\, \max_\omega\, \E[\Delta_T(\hat f, \by_\omega)]
    \;\ge\; \sum_{j=1}^k \delta_j \cdot \tfrac{1}{2}(1 - \TV_j)
    \;\ge\; \sum_{j=1}^k \frac{M^2 s_j}{2N}\big(1 - \tfrac{n s_j}{N}\big).
\end{equation*}
    Optimise per-block: $s_j^* := N/(2n)$, giving
    $\delta_j(1 - \TV_j) \le M^2/(8n) \cdot (1 - 1/2) = M^2/(16n)$.
    Summed over $k \le N/s_j^* = 2n$ blocks:
\begin{equation*}
    \inf_{\hat f}\, \max_\omega\, \E[\Delta_T] \;\le\; \frac{M^2 \cdot 2n}{16n} \;=\; \frac{M^2}{8}.
\end{equation*}
    This is $O(1)$, independent of $N - n$, confirming
    \eqref{eq:assouad-fails}. \qed
\end{enumerate}
\end{proof}

\begin{remark}[The structural reason]
\label{rem:assouad-why-fails}
The $1/\sqrt{N-n}$ Berry--Esseen rate of
\texttt{notes/04}~Theorem~3 originates from \emph{anticoncentration
of the empirical mean of a $[-M, M]$-valued function on a
without-replacement sample}, which is a $\sqrt{N-n}$-scale Gaussian
fluctuation. Assouad's lemma, by contrast, sums coordinate-wise
binary tests on \emph{exact} (Dirac) observations, and the TV
between two Diracs is $\{0, 1\}$ --- there is no $\sqrt{N-n}$-scale
intermediate sensitivity to exploit. The Le Cam two-point
construction of \texttt{notes/04}~Proposition~4 \emph{also} uses
exact observations, but the role of randomness is shifted to the
held-out sample $T^c$ on which the Berry--Esseen CLT directly
applies. Le Cam therefore captures the right rate; Assouad does not.
\end{remark}

\subsection{Combined statement: the operational gap is $\approx 3$}

\begin{theorem}[Sharp minimax rate and constant, v2]
\label{thm:axis-3-final}
For $n/N$ bounded away from $0, 1$ and $N - n \to \infty$, and
$\eta < 1/3$, the minimax slack rate of the leave-node-out excess
risk for any learner satisfying \texttt{paper.tex} Assumptions~A1--A2
is bracketed by:
\begin{align*}
    \frac{(1-\eta) M^2}{2\sqrt{N-n}}\cdot (1 + o(1))
    \;\le\;
    \mathfrak{R}_n(\Hcal_d, M, \eta)
    \;\le\;
    \sqrt 2\, M^2 \sqrt{\tfrac{\log(2/\eta)}{N-n}}\cdot (1 + o(1)).
\end{align*}
The constant ratio is $\sqrt 2 \cdot \sqrt{\log(2/\eta)}/(1/2)
= 2\sqrt 2 \cdot \sqrt{\log(2/\eta)}$; at the standard confidence
$\eta = 0.1$, $\sqrt{\log 20} \approx 1.73$, the gap is
$2\sqrt 2 \cdot 1.73 \approx 4.9$. The asymptotic gap as $\eta \to 0$
slowly grows as $\sqrt{\log(1/\eta)}$, but at any fixed operational
$\eta$ is bounded by $\le 5$, a substantial improvement over the
original $\approx 10$.
\end{theorem}

\begin{remark}[Why the factor-4 conjecture fails]
\label{rem:factor-4-fails}
The original \texttt{paper.tex}~\S8.3 item~(1) conjectured a factor
$\approx 4$ via multi-point Fano. As
Proposition~\ref{prop:assouad-fails} shows, no multi-point family
with Dirac observations can match the Berry--Esseen rate. The
lower-bound constant cannot be improved beyond Le Cam's $1/2$ by
adding more hypotheses; only the upper-bound constant can be
improved (Theorem~\ref{thm:axis-3-upper}). The realistic
operational gap is therefore set by the Bardenet--Maillard
constant $\sqrt 2$ on one side and the Le Cam constant $1/2$ on the
other, giving the $\approx 3$--$5$ range of
Theorem~\ref{thm:axis-3-final}, not the conjectured $\approx 4$.
\end{remark}

% =========================================================================
\section{Axis 4 --- Implicit-bias effective subspace}
\label{sec:axis-4}

\subsection{Setup: $\Omega$-tie-broken learner}

For an over-parameterised learner ($p \gg N$, e.g., MLP, kernel
ridgeless interpolation), the nominal hypothesis class $\Hcal_d$ may
satisfy $\Scal_\Acal = \R^N$, making Theorem~1 trivial
($\Rspec = 1$). Yet empirically, the optimiser's implicit bias
restricts the realised solutions to a much smaller subset
$\Scal_\Acal^\Omega$ of $\Hcal_d$. We formalise this:

\begin{definition}[$\Omega$-tie-broken learner]
\label{def:omega-learner}
Let $\Omega : \Hcal_d \to \R \cup \{+\infty\}$ be convex and
lower-semicontinuous. The \emph{$\Omega$-tie-broken ERM} is
\begin{equation}
    \hat f_\Acal^\Omega(T, \by') \;:=\;
    \argmin\!\big\{\Omega(f) : f \in \argmin_{g \in \Hcal_d} L_T(g, \by')\big\}.
    \label{eq:omega-learner}
\end{equation}
The \emph{$\Omega$-effective reachable subspace} is the closure of
the learner's image:
\begin{equation}
    \Scal_\Acal^\Omega \;:=\; \mathrm{closure}\!\left\{\,
        \hat f_\Acal^\Omega(T, \by') \;:\; T \in \binom{V}{n},\;
        \by' \in \R^n\,\right\} \;\subseteq\; \Hcal_d.
    \label{eq:s-omega}
\end{equation}
\end{definition}

\subsection{Universal bound on the effective subspace}

\begin{theorem}[Theorem~1 applies verbatim to $\Scal_\Acal^\Omega$]
\label{thm:axis-4}
Suppose $\Acal$ is the $\Omega$-tie-broken ERM
(Def.~\ref{def:omega-learner}). Then Theorem~1 of \texttt{notes/01}
holds with $\Scal_\Acal$ replaced by $\Scal_\Acal^\Omega$:
\begin{equation}
    \E_T\!\left[L_{T^c}(\hat f_\Acal^\Omega(T))\right]
    \;\ge\; \big(1 - \Rspec(\Scal_\Acal^\Omega, \by)\big)\Var(\by),
    \label{eq:axis-4-main}
\end{equation}
provided $\bP_{\Scal_\Acal^\Omega}\by \in \Scal_\Acal^\Omega$
(which holds whenever $\Scal_\Acal^\Omega$ is a subspace, by
closure under orthogonal projection).
\end{theorem}

\begin{proof}
The proof of \texttt{notes/01}~Theorem~1 only uses A1--A4 through
two facts about the learner:
(i)~$\hat f_\Acal(T) \in \Scal_\Acal$ (Bessel step), and
(ii)~$L_T(\hat f_\Acal(T)) \le L_T(\bP_{\Scal_\Acal}\by)$
(ERM step in expectation $= A$).

Substituting $\Scal_\Acal \to \Scal_\Acal^\Omega$:
(i) By construction $\hat f_\Acal^\Omega(T) \in \Scal_\Acal^\Omega$;
Bessel applies with the smaller subspace.
(ii) Since $\bP_{\Scal_\Acal^\Omega}\by \in \Scal_\Acal^\Omega$ and is
the minimiser of $\|f - \by\|^2$ on $\Scal_\Acal^\Omega$, it is also
the minimiser of $\E_T[L_T(f)]$; the $\Omega$-tie-broken ERM
satisfies the ERM step on $\Scal_\Acal^\Omega$ since $\inf$ over
$\Hcal_d$ equals $\inf$ over $\Scal_\Acal^\Omega$ for the squared
loss (the $\Omega$-tie-breaker selects a specific minimiser without
changing the infimum). Hence $\E_T[L_T(\hat f_\Acal^\Omega)] \le
\E_T[L_T(\bP_{\Scal_\Acal^\Omega}\by)] = (1 - \Rspec(\Scal_\Acal^\Omega, \by))\Var(\by)$.

The transduction identity is independent of the hypothesis class.
\qed
\end{proof}

\subsection{The operational difficulty for PD kernels}

\begin{remark}[When Theorem~\ref{thm:axis-4} is operationally
trivial]
\label{rem:axis-4-trivial}
Theorem~\ref{thm:axis-4} is mathematically correct but
operationally vacuous in the most important regime --- the
NTK / kernel-ridgeless setting where the optimiser is
SGD-with-small-init and the implicit bias is the RKHS norm.
By the representer theorem, the interpolant lives in
$\mathrm{range}(\bK) \subseteq \R^N$, but for strictly
positive-definite kernels (Gaussian RBF, NTK with non-degenerate
data) $\mathrm{range}(\bK) = \R^N$, so $\Rspec(\Scal_\Acal^\Omega, \by) = 1$
and the bound \eqref{eq:axis-4-main} is trivial.

The non-triviality must come from a finer characterisation of the
$\Omega$-tie-broken learner's image: not just the \emph{range} of the
kernel, but the \emph{minimum-norm interpolation map's effective
output subspace}, controlled by the spectral decay of $\bK$. This
requires a benign-overfitting-style analysis along the lines of
\citet{BartlettLongLugosiTsigler2020} adapted to the transductive
setting --- bounding the held-out excess risk as a function of the
kernel's effective rank, the noise level, and the target's
alignment with the top eigenvectors.

The full development of this bound is an open problem. The
research programme of
\texttt{notes/08\_benign\_overfitting\_research\_programme.tex}
formulates it as a precise sequence of conjectures
(Conjectures~1--3 there): under leverage-incoherence and
spectral-decay hypotheses, the held-out excess risk is
\emph{conjectured} to satisfy
$\E[L_{T^c}] \le C_1\tau_{r^*}\Var(\by) + C_2 r^*\mu^*/(n - r^*)
\Var(\by) + C_3\sigma_\xi^2 r^*/n + C_4 \sigma_\xi^2 n/\Reff_{r^*}(\bK)$,
where $(\reff, \Reff)$ are the BLLT effective rank/dimension and
$\mu^*$ is the leverage incoherence. A previous draft of
\texttt{notes/08} (now retracted) claimed a proof of this bound by
direct adaptation of BLLT, but the adaptation requires three
non-trivial new technical tools (matrix Chernoff for sampled
principal submatrices, exchangeable-martingale concentration for
$T$-dependent trace functionals, and a correct RKHS-vs-$\ell_2$
norm conversion), none of which were established in v1. The
present version of \texttt{notes/08} catalogues these obstacles
explicitly and points to the relevant literature
(\citealp{Tropp2012,Cohen2016sparsification,DerezinskiLiaoDoblerMahoney2020,Chatterjee2007exchangeable})
for the tools needed to make the conjectures into theorems.
\end{remark}

\begin{remark}[Non-$\ell_2$ implicit biases]
\label{rem:lasso}
For $\Omega = \|\cdot\|_1$ on a coordinate basis (Lasso, hard
thresholding, ReLU-NN with sparse init), $\Scal_\Acal^\Omega$ is
\emph{not} a subspace but a union of $\le \binom{p}{n}$ coordinate
faces of dimension $\le n$. The bound \eqref{eq:axis-4-main}
generalises by replacing $\Rspec$ with the supremum:
\begin{equation*}
    \Rspec^{\Omega}(\by)
    \;:=\; \sup_{S \subset [p],\, |S| \le n} \Rspec(\Scal_S, \by),
\end{equation*}
where $\Scal_S$ is the coordinate subspace indexed by $S$. The
proof is identical with $\Scal_\Acal^\Omega \subseteq \bigcup_S
\Scal_S$ replacing the subspace inclusion in step~(i). A sharp
characterisation of $\Rspec^\Omega$ for adversarially-chosen sparse
supports is an open question.
\end{remark}

% =========================================================================
\section{Synthesis and combined statement}
\label{sec:synthesis}

The four axes are largely independent on the proof level.
Axes~1 and 2 admit a clean combined statement under linear ERM;
Axis~3 sharpens the constant on both upper-bound and lower-bound
sides of the minimax rate; Axis~4 is a reparametrisation of the
hypothesis-class object on which all three previous axes operate.

\subsection{Implications for \texttt{paper.tex}}
\label{sec:paper-implication}

The required edits to \texttt{paper.tex}~\S8 are:
\begin{itemize}[label=$\bullet$,leftmargin=*]
\item \textbf{\S8.3 item~(1) (sharp constant):} replace
    ``factor $\approx 4$'' by ``factor $\approx 3$--$5$ at typical
    operational confidence $\eta = 0.1$''; cite
    \citet{BardenetMaillard2015} for the upper-bound refinement
    (Theorem~\ref{thm:axis-3-upper}); cite
    Proposition~\ref{prop:assouad-fails} for the negative result
    on Assouad. The multi-point Fano route is closed.
\item \textbf{\S8.3 item~(3) (approximate ERM):} the slack
    $(\rho-1)\E_T[\epsilon_T]$ is now rigorously derived
    (Theorem~\ref{thm:axis-2}), with calibration table for ridge /
    GBM / NN.
\item \textbf{\S8.3 item~(4) (misspecified $\Hcal_d$):} the bound
    applies verbatim to the $\Omega$-implicit-bias subspace
    $\Scal_\Acal^\Omega$ (Theorem~\ref{thm:axis-4}); however, for
    strictly-PD NTK-regime kernels this remains operationally
    trivial ($\Rspec = 1$). A non-trivial floor in this regime is
    \emph{an open problem} catalogued as a research programme in
    \texttt{notes/08\_benign\_overfitting\_research\_programme.tex},
    not a result. The original Mercer-truncation route via simple
    Pythagore (an earlier draft of this note) was incorrect:
    $\hat f \in \mathrm{range}(\bK)$ does not project cleanly onto
    $\Scal_r$.
\item \textbf{\S8.4 Limitations (deterministic target):} the
    stochastic extension is now Theorem~\ref{thm:axis-1}
    (linear ERM, distribution-free noise) and
    Theorem~\ref{thm:axis-1-stein} (non-linear ERM, Gaussian noise
    via Stein--SURE), with slack $-2 d_{\mathrm{eff}}\sigma_\xi^2/(N-n)$
    matching the in-sample noise variance of OLS at fixed design.
    A data-independent-hyperparameter assumption (A\ref{a:dataindep})
    is explicit.
\end{itemize}

\subsection{What remains genuinely open}

\begin{itemize}[label=$\bullet$,leftmargin=*]
\item \textbf{Constant gap $< 3$.} Achieving constant
    ratio $< 3$ requires either (a) a constant
    $c_{\mathrm{LB}} > 1/2$ on the lower bound (impossible via
    Assouad/Fano, may be possible via local-asymptotic-minimax
    \`a la \citet{LeCam1986}~Ch.~6), or (b) a Rademacher constant
    $< \sqrt 2$ on the upper bound (would require beyond
    Bardenet--Maillard; permutation-Rademacher of
    \citet{Pechyony2008phd}~Thm.~3.4 gives $\sqrt{\pi/2} \approx 1.25$
    but requires extra hypotheses on the loss).
\item \textbf{NTK-regime tight floor.} The transductive
    benign-overfitting programme of \texttt{notes/08} formulates
    the bound as Conjectures~1--3 parameterised by missing tools
    (matrix Chernoff for sampled principal submatrices,
    exchangeable-martingale concentration for $T$-dependent
    trace functionals, RKHS-vs-$\ell_2$ norm conversion). All three
    are within reach of established random-matrix and exchangeable
    concentration literature, but writing them down in the
    leave-node-out setting is a research programme of its own.
\item \textbf{CV-driven hyperparameter selection (Axis~1).} The
    stability/PAC-Bayes correction to A\ref{a:dataindep} for
    CV-tuned $\lambda$ is computable but not yet derived in
    closed form for ridge.
\item \textbf{Combined noisy-and-approximate ERM.} The slacks of
    Axes~1 and 2 combine via the cross-term in
    \eqref{eq:axis-1-erm-noisy}, not via a simple sum or max
    (Remark~\ref{rem:axis-12-combined}). Derivation of the exact
    combined slack is mechanical but tedious.
\item \textbf{Non-$\ell_2$ implicit biases.} The Lasso analogue
    (Remark~\ref{rem:lasso}) needs a sharp characterisation of
    $\sup_S \Rspec(\Scal_S, \by)$ for adversarially-chosen sparse
    supports.
\item \textbf{Importance-weighted transduction} for active-learning
    protocols: \texttt{paper.tex}~\S8.3~item~(2), deferred to the
    C3 companion paper.
\end{itemize}

% =========================================================================
\bibliography{../references/references}

\end{document}
